\documentclass[]{article}

\usepackage{amsmath}
\usepackage{multirow}
\usepackage{natbib}
\usepackage{graphicx} 


\begin{document}

The study of anomalous transport, where particle diffusion deviates from the classical Brownian paradigm, is a central theme in statistical physics. In contrast to normal diffusion, where the mean squared displacement (MSD) grows linearly with time, $\langle \mathbf{r}^2(t) \rangle \sim t$, many systems in nature exhibit superdiffusion, characterized by an MSD scaling of $\langle \mathbf{r}^2(t) \rangle \sim t^\alpha$ with an exponent $\alpha > 1$. While such behavior is often modeled using stochastic Lévy processes, which are defined by random jumps drawn from heavy-tailed probability distributions, a fundamental question remains: how do these statistical frameworks apply to deterministic systems where anomalous transport arises not from temporal randomness, but from the complex, spatially correlated structure of a quenched environment?

A canonical system for exploring this question is the advection of passive tracers in a velocity field generated by a collection of vortices. In two dimensions, the flow induced by point vortices is well-known to generate superdiffusion governed by Cauchy-Lévy statistics. The three-dimensional counterpart, a velocity field generated by a static arrangement of vortex filaments, presents a distinct and less-explored problem. The velocity field in this case is dictated by the Biot-Savart law, which for a single filament induces a velocity that decays with distance $r$ as $1/r$. For a random collection of filaments, the superposition of these fields is theoretically predicted to produce a heavy-tailed velocity probability distribution known as the Holtsmark distribution. This prediction suggests a unique universality class for anomalous transport, corresponding to a Lévy walk with an expected MSD exponent of $\alpha = 1.5$. However, it is unclear how these asymptotic statistical predictions manifest in a finite-density system and what physical mechanisms connect the flow's deterministic geometry to the emergent transport statistics.

In this work, we address this challenge through direct numerical simulations of tracer transport in a three-dimensional quenched velocity field generated by static vortex filaments. We systematically characterize the transport regime across a range of filament densities to investigate the convergence towards the theoretically predicted asymptotic limit. By analyzing tracer trajectories, we directly measure the anomalous diffusion exponent from the MSD and compare the empirical velocity probability distribution with the theoretical Holtsmark distribution. Going beyond a purely statistical description, we establish a mechanistic link between the transport dynamics and the local topology of the flow field, demonstrating that low-speed trapping events are localized in rotation-dominated regions. Furthermore, we analyze velocity correlations to quantify the persistent memory inherent in this deterministic system, highlighting a crucial distinction from memoryless stochastic processes like canonical Lévy walks. Our results provide a comprehensive picture of how quenched spatial disorder in a three-dimensional vortex field shapes a distinct, non-ergodic superdiffusive regime, bridging the gap between abstract statistical theory and concrete physical dynamics.

\end{document}
                